\begin{abstract}
Examination boards publish results at an announced time, and their websites
often fail under the rush of students. We built ResultBench, a testbed that
runs eleven designs of a result portal under the same two-CPU budget, replays
a result-day surge of students who wait and refresh, and checks every answer
against ground truth. On one machine, a database-backed portal served every
student within \SI{5}{s} at its calibrated capacity but only 30.8\% at twice
that load, and neither mirror sites nor an empty cache helped. Pre-computing
results into Redis or static files served every student at ten times the load
with 12.5--57.8 times less CPU per student, after less than \SI{10}{s} of
preparation. No design returned wrong results; stale results after
corrections came only from cached copies that were not updated. Because every
student asks for a different result, almost every cache miss was the first
request for its key, so stampede defences changed nothing, while filling the
cache before release served everyone. Refreshing students almost doubled the
requests per student, and impatient ones kept an unprotected portal
overloaded. A virtual waiting room at five times the load turned failures into
long waits rather than serving more students. We make ResultBench available so
that boards can test their own designs.
\end{abstract}
